\documentclass[11pt]{article}
\usepackage[a4paper,margin=25mm]{geometry}
\usepackage{amsmath,amssymb,amsthm}
\usepackage[hidelinks]{hyperref}
\emergencystretch=2em
\newtheorem{theorem}{Theorem}
\newtheorem{lemma}{Lemma}
\newcommand{\per}{\operatorname{per}}
\title{Exact sign pairing in the permanent formula\\
\large Proof of Conjecture 00000006322}
\author{gaochengzhecpu}\date{}
\begin{document}\maketitle
\begin{abstract}
For an $n\times n$ matrix with $n\ge1$, we derive the full
$2^n$-term signed product-of-sums formula for the permanent and
prove that its terms pair under simultaneous sign reversal.
Each pair contributes twice either representative. Fixing one
coordinate to $+1$ leaves exactly $2^{n-1}$ terms, with normalization
$2^{-(n-1)}$. The Lean proof covers every positive matrix size and
every field in which $2\ne0$, and uses Mathlib's standard permanent.
\end{abstract}

\section*{Precise statement and conventions}
Let $I$ be a finite set of size $n\ge1$, let $R$ be a field with
$2\ne0$, and let $A=(a_{ij})_{i,j\in I}$ be a square matrix over $R$.
In particular the result applies over the rationals, reals or
complex numbers. Define
\[
 \per(A)=\sum_{\sigma\in\operatorname{Sym}(I)}
                 \prod_{j\in I}a_{\sigma(j),j}.
\]
For a sign vector $\delta\in\{1,-1\}^I$ put
\[
 W_A(\delta)=\left(\prod_{i\in I}\delta_i\right)
                 \prod_{j\in I}\left(\sum_{i\in I}\delta_i a_{ij}\right).
\]
Choose any coordinate $i_0\in I$. We prove the following exact
form of the source's pairing and conversion assertions.

\begin{theorem}\label{thm:main}
For every such matrix,
\[
 \boxed{\displaystyle
 \per(A)=\frac{1}{2^{n-1}}
       \sum_{\substack{\delta\in\{1,-1\}^I\\\delta_{i_0}=1}}W_A(\delta).}
\]
The displayed sum has exactly $2^{n-1}$ indexed terms. It is obtained
from the full $2^n$-term formula by the fixed-point-free pairing
$\delta\leftrightarrow-\delta$, and paired terms are equal.
\end{theorem}

This is the classical sign-sum permanent identity commonly called
Glynn's formula; a published elementary derivation is given by
Spies. We supply the full proof and a general Lean formalization.
The term count describes this particular representation, not a
lower bound on all algorithms or a claim that every summand is
nonzero for every matrix. Positive size is required for the free
pairing. The empty permanent has its separate convention $1$.

\section*{Derivation of the full formula}
Distributing the $n$ column sums gives
\[
 W_A(\delta)=\sum_{f:I\to I}C_f(\delta)
                     \prod_{j\in I}a_{f(j),j},\qquad
 C_f(\delta)=\left(\prod_{i\in I}\delta_i\right)
                    \left(\prod_{j\in I}\delta_{f(j)}\right).
\]
We evaluate the sum of $C_f$ over all sign vectors.

If $f$ is bijective, then $\prod_j\delta_{f(j)}=\prod_i\delta_i$.
Therefore $C_f(\delta)=1$ for every sign vector, since
$\delta_i^2=1$. Its sum is $2^n$.

If $f$ is not bijective, finiteness of its common domain and codomain
implies that $f$ is not surjective. Choose $r\notin f(I)$.
Flipping only $\delta_r$ negates the first product in $C_f$ and
leaves the second unchanged. This flip is an involutive bijection
of the sign vectors. Thus its sum $S_f$ satisfies $S_f=-S_f$;
since $2\ne0$ in a field, $S_f=0$.

Interchanging the two finite sums consequently yields
\begin{align*}
 \sum_{\delta\in\{1,-1\}^I}W_A(\delta)
 &=\sum_{f:I\to I}\left(\sum_\delta C_f(\delta)\right)
                           \prod_{j\in I}a_{f(j),j}\\
 &=2^n\sum_{\sigma\in\operatorname{Sym}(I)}
                           \prod_{j\in I}a_{\sigma(j),j}
 =2^n\per(A).
\end{align*}
This proves the full formula directly from the permutation
definition; the full formula is not an additional assumption.

\section*{The pairing and its exact constant}
Simultaneously negating all coordinates contributes $(-1)^n$
from the prefactor and another $(-1)^n$ from the $n$ column sums.
Consequently
\[
 W_A(-\delta)=(-1)^{2n}W_A(\delta)=W_A(\delta).
\]
Because $n\ge1$ and $1\ne-1$, no sign vector equals its negative.
Every two-element pair has exactly one member with
$\delta_{i_0}=1$. These representatives are indexed freely by
the other $n-1$ signs, so there are exactly $2^{n-1}$ of them.
It follows that
\[
 2^n\per(A)=\sum_\delta W_A(\delta)
       =2\sum_{\delta_{i_0}=1}W_A(\delta).
\]
Cancel $2$ and then divide by the nonzero $2^{n-1}$ to obtain
Theorem~\ref{thm:main}. Thus the normalization changes from
$2^{-n}$ to $2^{-(n-1)}$, exactly compensating for retaining one
representative of each equal pair. This proves all three pairing,
constant and term-count assertions in the stated formula.

\section*{Formal correspondence}
The Lean result is generic over a finite index type, a field $R$,
and the assumption \texttt{NeZero (2 : R)}. The entries of the
matrix are arbitrary. \texttt{Matrix.permanent} is used directly.
Signs are encoded by Booleans: false is $+1$ and true is $-1$.
The theorem \texttt{sign\_injective} verifies that these are distinct
field values under the stated assumption.

The theorem \texttt{term\_expansion} distributes the column sums;
\texttt{coefficient\_sum} proves the cancellation for every
nonbijective row-choice function. An explicit equivalence between
bijective functions and permutations connects the surviving terms
to the actual permanent in \texttt{full\_formula}.

For the halving step, \texttt{opposite\_involutive} and
\texttt{opposite\_ne\_self} establish the free pairing.
\texttt{term\_opposite} proves equality of its summands, and
\texttt{half\_sums\_equal} constructs the bijection between the two
halves. \texttt{number\_of\_terms} proves the cardinality
$2^{|I|-1}$. The final theorem \texttt{pairing\_reduction} combines
the free pairing, the factor two, the exact count and the permanent
identity. \texttt{glynn\_succ} specializes it to every matrix of
size $n+1$, including the one-by-one case.

\section*{Reproduction}
Lean 4.19.0 and Mathlib revision\\
\texttt{c44e0c8ee63ca166450922a373c7409c5d26b00b}
are pinned. In \texttt{lean/}, run
\begin{verbatim}
lake exe cache get
lake build
lake env lean -DwarningAsError=true Main.lean
\end{verbatim}
The submitted project is rebuilt in a fresh directory using only
verified official dependency artifacts. No numerical test family
substitutes for the general proof. There are no proof placeholders,
custom axioms or external computational oracles. No auxiliary
numerical program is required. Run \texttt{tectonic main.tex} to
regenerate the PDF. The author performed a separate adversarial
self-review; no independent reviewer or subagent was used.

\section*{References}
\begin{enumerate}
\item The Last Math Competition, Conjecture 00000006322. The exact
bilingual source and its upstream provenance are included.
\item J. Spies, \emph{A formula for the permanent}, Nieuw Archief
voor Wiskunde 5/21(1) (2020), p.~27.
\href{https://www.nieuwarchief.nl/serie5/pdf/naw5-2020-21-1-027.pdf}{Publisher's article}.
\end{enumerate}
\end{document}
